\begin{figure*}[t]
\centering
\begin{tikzpicture}
\begin{axis}[
    width=0.86\textwidth,
    height=0.40\textwidth,
    xlabel={恢复的自然缺口（共 3,155 个）},
    ylabel={实测失败比例（\%）},
    xmin=315, xmax=545,
    ymin=12, ymax=64,
    grid=major,
    grid style={gray!18},
    legend pos=north east,
    legend cell align=left,
    legend style={draw=none, fill=none, font=\small},
    tick label style={font=\small},
    label style={font=\small},
    mark size=2.4pt,
]
\addplot[only marks, mark=o, mark options={solid}, color=blue!65!black,
         forget plot] coordinates {
    (529,59.95) (399,57.73) (388,46.77) (391,44.30) (395,48.05)
    (393,35.39) (393,47.15) (383,21.15) (383,21.24)
};
\addlegendimage{only marks, mark=o, color=blue!65!black}
\addlegendentry{ByteTrack}
\addplot[only marks, mark=square, mark options={solid}, color=red!65!black,
         forget plot] coordinates {
    (424,53.39) (354,48.69) (343,37.76) (347,41.02) (339,25.19)
    (346,37.28) (344,37.78) (327,18.31) (334,16.19)
};
\addlegendimage{only marks, mark=square, color=red!65!black}
\addlegendentry{BoT-SORT（仅运动）}
\addplot[only marks, mark=star, mark size=5pt, color=black] coordinates {(529,59.95)};
\addplot[only marks, mark=star, mark size=5pt, color=black] coordinates {(424,53.39)};
\addplot[only marks, mark=triangle*, mark size=4pt, color=blue!65!black] coordinates {(393,35.39)};
\addplot[only marks, mark=triangle*, mark size=4pt, color=red!65!black] coordinates {(334,16.19)};
\node[font=\scriptsize, anchor=south west] at (axis cs:531,59.0) {不限额};
\node[font=\scriptsize, anchor=north east] at (axis cs:391,34.2) {交互};
\node[font=\scriptsize, anchor=south east] at (axis cs:332,16.9) {严格};
\end{axis}
\end{tikzpicture}
\caption{官方 PersonPath22 测试子集上的全部已评估策略，横轴为恢复的自然缺口数，
纵轴为可评估重连中的实测失败比例。星号标出不限额跟踪器，三角标出松弛交互门控，
左下角一组为严格天花板。靠右的策略恢复更多缺口，靠下的策略更少失败。前沿从左下的
严格工作点延伸到右上的不限额原版，没有任何一个策略在两个轴上都最优。}
\label{fig:tradeoff-zh}
\end{figure*}
